\KOMAoption{paper}{a4, portrait}
\areaset{267mm}{180mm}
\fancyheadoffset{0pt} % обновляем хедер и футер

\newgeometry{
  top=14mm,
  bottom=11mm,
  right=10mm,
  left=22mm,
}

\begin{landscape}

\setlength{\footskip}{15.5pt}

\color{uniblue}{\section[(рекомендуемое) Схемы подключения цепей треугольника]{(рекомендуемое)\\Схемы подключения цепей треугольника}\label{app:analogs}}
\color{black}

\captionsetup[table]{margin=0.6cm} % временное изменение
\small
\setlength{\extrarowheight}{0.05cm} % высота ячеек в таблице

\begin{longtable}{
  >{\centering\arraybackslash}m{0.5cm}
  |>{\centering\arraybackslash}m{5.5cm}
  |>{\centering\arraybackslash}m{1.0cm}
  |>{\centering\arraybackslash}m{3.5cm}
  |>{\centering\arraybackslash}m{5.54cm}
  |>{\centering\arraybackslash}m{5.4cm}
  |>{\centering\arraybackslash}m{3.0cm}
  }

\caption{Схемы подключения цепей треугольника дополнительной обмотки ТН: $«U_\textsl{НИ}/U_\textsl{ИК}»$\hfill\vspace{-0.5\baselineskip}}\label{desc_tznp:tbl}\\ 
\hhline{|-|-|-|-|-|-|-|}
\rowcolor{gray!30}
№ & Схема соединения обмоток «разомкнутого треугольника» ТН  & Особая фаза  & Направление векторов напряжений «звезды» и «треугольника»  & Векторная диаграмма напряжений  & Выражение для расчета $U_{\textsl{БНН}}$ & Выражение для расчета ($3U_0$ треугольник)   \\ 
\hhline{|-|-|-|-|-|-|-|}
\endfirsthead

\caption*{\hspace{3pt}\emph{Продолжение таблицы \ref{desc_tznp:tbl}\hfill\vspace{-0.5\baselineskip}}} \\ % сделано по ГОСТ 2.105 п.6.8.7
\hhline{|-|-|-|-|-|-|-|}
\rowcolor{gray!30}
№ & Схема соединения обмоток «разомкнутого треугольника» ТН  & Особая фаза  & Направление векторов напряжений «звезды» и «треугольника»  & Векторная диаграмма напряжений  & Выражение для расчета $U_{\textsl{БНН}}$ & Выражение для расчета ($3U_0$ треугольник)   \\ 
\hhline{|-|-|-|-|-|-|-|}
\endhead
\hhline{|-|-|-|-|-|-|-|}

\endfoot
\endlastfoot

1  & \includegraphics[width=0.3\textwidth,height=0.3\textheight,keepaspectratio]{\apppath/Приложение. БНН/img/111.pdf}  &  $A$  &  Совпадает &  \vspace{2mm}\includegraphics[width=0.3\textwidth]{\apppath/Приложение. БНН/img/112.pdf} &  $ \vec U_{\textsl{БНН}} = -\vec U_A + \vec U_B + \vec U_C + \frac{\vec U_\textsl{НИ} - \vec U_\textsl{ИК}}{\sqrt{3}} $  &  $U_\textsl{НИ}+U_\textsl{ИК}$ \\
\hline

2  & 
\includegraphics[width=0.3\textwidth,height=0.3\textheight,keepaspectratio]{\apppath/Приложение. БНН/img/121.pdf}  &  $A$  &  >> & \vspace{2mm}\includegraphics[width=0.3\textwidth]{\apppath/Приложение. БНН/img/122.pdf}  &  >> &  >> \\
\hline

3  & 
\includegraphics[width=0.3\textwidth,height=0.3\textheight,keepaspectratio]{\apppath/Приложение. БНН/img/131.pdf}  &  $A$  &  Не совпадает &  \vspace{2mm}\includegraphics[width=0.3\textwidth]{\apppath/Приложение. БНН/img/132.pdf} &  $ \vec U_{\textsl{БНН}} = -\vec U_A + \vec U_B + \vec U_C - \frac{\vec U_\textsl{НИ} - \vec U_\textsl{ИК}}{\sqrt{3}} $  &  $U_\textsl{НИ}+U_\textsl{ИК}$ \\
\hline

4  & \includegraphics[width=0.3\textwidth,height=0.3\textheight,keepaspectratio]{\apppath/Приложение. БНН/img/141.pdf}  &  $A$  &  >> & \vspace{1mm}\includegraphics[width=0.3\textwidth]{\apppath/Приложение. БНН/img/142.pdf}  &  >> &  >> \\
\hline

5  & \includegraphics[width=0.3\textwidth,height=0.3\textheight,keepaspectratio]{\apppath/Приложение. БНН/img/151.pdf}  &  $B$  &  Совпадает &  \vspace{1mm}\includegraphics[width=0.3\textwidth]{\apppath/Приложение. БНН/img/152.pdf} &  $ \vec U_{\textsl{БНН}} = \vec U_A - \vec U_B + \vec U_C + \frac{\vec U_\textsl{НИ}- \vec U_\textsl{ИК}}{\sqrt{3}} $  &  $U_\textsl{НИ}+U_\textsl{ИК}$ \\
\hline

6  & \includegraphics[width=0.3\textwidth,height=0.3\textheight,keepaspectratio]{\apppath/Приложение. БНН/img/161.pdf}  &  $B$  &  >> & \vspace{0mm}\includegraphics[width=0.3\textwidth]{\apppath/Приложение. БНН/img/162.pdf}  &  >> &  >> \\
\hline

7  & \includegraphics[width=0.3\textwidth,height=0.3\textheight,keepaspectratio]{\apppath/Приложение. БНН/img/171.pdf}  &  $B$  &  Не совпадает &  \vspace{0mm}\includegraphics[width=0.3\textwidth]{\apppath/Приложение. БНН/img/172.pdf} &  $ \vec U_{\textsl{БНН}} = \vec U_A - \vec U_B + \vec U_C - \frac{\vec U_\textsl{НИ} - \vec U_\textsl{ИК}}{\sqrt{3}} $  &  $U_\textsl{НИ}+U_\textsl{ИК}$ \\
\hline

8  & 
\includegraphics[width=0.3\textwidth,height=0.3\textheight,keepaspectratio]{\apppath/Приложение. БНН/img/181.pdf}  &  $B$  &  >> & \vspace{0mm}\includegraphics[width=0.3\textwidth]{\apppath/Приложение. БНН/img/182.pdf}  &  >> &  >> \\
\hline


9  & \includegraphics[width=0.3\textwidth,height=0.3\textheight,keepaspectratio]{\apppath/Приложение. БНН/img/191.pdf}  &  $C$  &  Совпадает &  \vspace{0mm}\includegraphics[width=0.3\textwidth]{\apppath/Приложение. БНН/img/192.pdf} &  $ \vec U_{\textsl{БНН}} = \vec U_A + \vec U_B- \vec U_C + \frac{\vec U_\textsl{НИ} - \vec U_\textsl{ИК}}{\sqrt{3}} $  &  $U_\textsl{НИ}+U_\textsl{ИК}$ \\
\hline

10  & 
\includegraphics[width=0.3\textwidth,height=0.3\textheight,keepaspectratio]{\apppath/Приложение. БНН/img/1101.pdf}  &  $C$  &  >> & \vspace{0mm}\includegraphics[width=0.3\textwidth]{\apppath/Приложение. БНН/img/1102.pdf}  &  >> &  >> \\
\hline

11  & \includegraphics[width=0.3\textwidth,height=0.3\textheight,keepaspectratio]{\apppath/Приложение. БНН/img/1111.pdf}  &  $C$  &  Не совпадает &  \vspace{0mm}\includegraphics[width=0.3\textwidth]{\apppath/Приложение. БНН/img/1112.pdf} &  $ \vec U_{\textsl{БНН}} = \vec U_A + \vec U_B - \vec U_C - \frac{\vec U_\textsl{НИ} - \vec U_\textsl{ИК}}{\sqrt{3}} $  &  $U_\textsl{НИ}+U_\textsl{ИК}$ \\
\hline

12  & \includegraphics[width=0.3\textwidth,height=0.3\textheight,keepaspectratio]{\apppath/Приложение. БНН/img/1121.pdf}  &  $C$  &  >> & \vspace{0mm}\includegraphics[width=0.3\textwidth]{\apppath/Приложение. БНН/img/1122.pdf}  &  >> &  >> \\
\hline











\end{longtable}
\setlength{\extrarowheight}{0.0cm} % возврат параметра высота ячеек в таблице
\normalsize
\captionsetup[table]{margin=0cm} % вернуть обратно если нужно


















\clearpage
\end{landscape}
\restoregeometry


